\documentclass[12pt]{article}
\usepackage[T1]{fontenc}
\usepackage[utf8]{inputenc}
\usepackage{lmodern}
\usepackage{textcomp}
\usepackage{enumitem}
\usepackage{geometry}
\usepackage{graphicx}
\usepackage{amsmath, amssymb}
\geometry{margin=1in}
\begin{document}
\begin{center}
Constitutional Law - Graduate Level Assessment 3 \
Total Marks: 40 \
Answer all questions.
\end{center}

\section*{Multiple Choice (10 marks)}
\begin{enumerate}[label=\arabic*.]
\item Which of these statements best describes the UK Constitution?
\begin{enumerate}[label=(\alph*)]
    \item The UK Constitution is based on a Bill of Rights
    \item The UK Constitution gives the judiciary the power to overturn acts of parliament
    \item The UK Constitution's power is derived solely from Parliament
    \item The UK Constitution is codified and contained in a single document
    \item The UK Constitution is uncodified and can be found in a number of sources
\end{enumerate}
\item Who said that "Jurisprudence is the eye of law"
\begin{enumerate}[label=(\alph*)]
    \item Pound
    \item Austin
    \item Savigny
    \item Hart
    \item Bentham
\end{enumerate}
\item What is the function of "standard-setting in human rights diplomacy?
\begin{enumerate}[label=(\alph*)]
    \item Standard-setting means setting certain standards of conduct in human rights treaties
    \item Standard-setting means putting forward both binding and non-binding legal standards
    \item Standard-setting means putting forward non-binding legal standards
    \item Standard-setting means enforcing human rights treaties
    \item Standard-setting means establishing a set of universal human rights principles
\end{enumerate}
\item What was the outcome before the European Court of Human Rights in the Al-Adsani case?
\begin{enumerate}[label=(\alph*)]
    \item The Court held that the right to a fair trial trumped the privilege of immunity
    \item The Court held that the privilege of immunity was not applicable in this case
    \item The Court held that the right to a fair trial was not applicable in this case
    \item The Court held that immunities were not in conflict with the right to a fair trial
    \item The Court held that the case was admissible due to overriding human rights considerations
\end{enumerate}
\item Which statement below best represents Durkheim's view of the function of punishment?
\begin{enumerate}[label=(\alph*)]
    \item Vengeance.
    \item Deterrence.
    \item Restoration.
    \item Moral education.
    \item Indemnification.
\end{enumerate}
\end{enumerate}

\section*{True / False (10 marks)}
\textit{Answer True or False}
\begin{enumerate}[label=\arabic*.]
\setcounter{enumi}{5}
\item True or False: Standard-setting in human rights diplomacy entails establishing binding legal standards that must be followed by states.
\item True or False: The principle of justice as fairness suggests that the original position would prioritize wealth over a compassionate society.
\item True or False: The defense power, standing alone, cannot support the enactment of federal law under the Constitution.
\item True or False: The case of Golder v UK (1978) is widely recognized as the leading example of the 'living instrument principle' by the European Court of Human Rights.
\item True or False: Critical legal theorists believe that law is a stable and predictable system that can be applied consistently.
\end{enumerate}

\section*{Long Form Response (20 marks)}
\textit{Answer each question with a clear and complete explanation.}
\begin{enumerate}[label=\arabic*.]
\setcounter{enumi}{10}
\item Discuss the International Court of Justice's stance on the accumulation theory regarding armed attacks, focusing on its implications for the interpretation of international law.
\item Discuss the admissibility of newspaper evidence regarding a corporation's stock price in a stock fraud case, analyzing its relevance and potential impact on the proceedings.
\end{enumerate}

\vfill
\noindent\textit{End of Paper}
\end{document}